\uuid{zfyp}
\exo7id{2647}
\titre{exo7 2647}
\auteur{debievre}
\organisation{exo7}
\datecreate{2009-05-19}
\isIndication{false}
\isCorrection{false}
\chapitre{Topologie}
\sousChapitre{Ouvert, fermé, intérieur, adhérence}
\cpge{mpsi-mp2i.fonctions-deux-variables.a, pcsi.fonctions-deux-variables.a, ptsi.fonctions-deux-variables.a, mp-mpi.topologie-espaces-normes.d, pc.espaces-vectoriels-normes.d, psi.espaces-vectoriels-normes.d, pt.fonctions-deux-trois-variables.a}
\module{Analyse}
\niveau{L2}
\difficulte{}

\contenu{
\texte{
Soit $D$ un sous-ensemble de $\R^n$.
}
\begin{enumerate}
    \item \question{Donner la d\'efinition de ``$D$ est ouvert.'' (Ceci est une question de cours!)}
    \item \question{Donner la d\'efinition de ``$a\in\R^2$ est un point adh\'erent de $D$.'' (Ceci est une question de cours!)

On consid\`ere dans la suite de l'exercice l'ensemble
$$D=\{(x,y)\in\R^2 \mid |x|\leq |y|, x^2+y^2<1\}\subset \R^2.$$}
    \item \question{Dessiner $D$.}
    \item \question{Montrer que $D$ n'est pas ouvert.}
    \item \question{D\'eterminer $\overline D$, l'adh\'erence de $D$. On justifiera bri\`evement sa r\'eponse, en s'aidant d'un dessin.}
\end{enumerate}
}
